\chapter{Mô tả bài toán}

Sudoku là trò chơi điền số trên lưới \(9 \times 9\). Lưới chia thành 9 khối \(3 \times 3\) bằng các đường kẻ đậm. Một lưới Sudoku hợp lệ khi mỗi hàng, mỗi cột và mỗi khối đều chứa đủ 9 số từ 1 đến 9, mỗi số đúng một lần.

Cho một lưới trong đó một số ô đã điền sẵn, các ô còn lại để trống. Hãy điền số vào các ô trống để được một lưới Sudoku hợp lệ, giữ nguyên các số đã cho.

\subsection*{Input}

Gồm 9 dòng, mỗi dòng 9 ký tự cách nhau bởi dấu cách. Mỗi ký tự là một số từ 1 đến 9 (ô đã điền) hoặc chữ \texttt{X} (ô trống). Số ô trống từ 0 đến 81.

\subsection*{Output}

\begin{itemize}
  \item Nếu điền được: 9 dòng, mỗi dòng 9 số cách nhau bởi dấu cách, là lưới đã điền kín và hợp lệ. Nếu có nhiều cách điền, xuất cách đầu tiên tìm được khi xét ô trống theo thứ tự hàng rồi cột và thử số từ nhỏ đến lớn.
  \item Nếu đề có hai ô cùng số trong cùng một hàng, cột hoặc khối: thông báo đề không hợp lệ, kèm số trùng và vị trí.
  \item Nếu không có cách điền nào: thông báo đề không có lời giải.
\end{itemize}

\subsection*{Ví dụ}

\begin{center}
  \begin{minipage}[t]{0.44\linewidth}
    \textbf{Input}
    \begin{Verbatim}[fontsize=\small, frame=single]
5 3 4 6 7 X 9 X X
6 7 X X 9 5 3 4 8
X 9 8 3 4 X 5 6 7
8 X 9 7 6 X X 2 3
4 X 6 X 5 3 7 9 1
7 1 3 9 X 4 8 5 6
9 6 1 5 3 7 2 X 4
2 8 7 4 1 9 6 3 5
X 4 5 X 8 6 1 7 9
\end{Verbatim}

  \end{minipage}
  \hfill
  \begin{minipage}[t]{0.44\linewidth}
    \textbf{Output}
    \begin{Verbatim}[fontsize=\small, frame=single]
5 3 4 6 7 8 9 1 2
6 7 2 1 9 5 3 4 8
1 9 8 3 4 2 5 6 7
8 5 9 7 6 1 4 2 3
4 2 6 8 5 3 7 9 1
7 1 3 9 2 4 8 5 6
9 6 1 5 3 7 2 8 4
2 8 7 4 1 9 6 3 5
3 4 5 2 8 6 1 7 9
\end{Verbatim}

  \end{minipage}
\end{center}

\subsection*{Yêu cầu}

Giải bằng kỹ thuật quay lui theo mô hình 1. Chương trình phải cho người dùng nhập đề, theo dõi từng bước thuật toán thực hiện, và xem kết quả.
